\section{A classical baseline and the cost}\label{sec:cost}

\subsection{The Ritz projection}\label{sec:cost:ritz}

The Ritz projection of \cref{sec:methods:ritz} needs the instance's stiffness matrix but no
labels or training (its rows in \cref{tab:3d,tab:2d}). On the 256 three-dimensional
validation instances its displacement error was smaller than the label-free transformer's on
\numRitzValFreeDispSixMin{} to \numRitzValFreeDispSixMax{} of the instances at $p=6$,
depending on the seed, and on \numRitzValFreeDispNineMin{} to \numRitzValFreeDispNineMax{}
at $p=9$; at $p=9$ it was smaller than each network's, in each seed, on at least
\numRitzValNineDispAll{} of the instances. Its median relative energy gap fell from
\numRitzValGapSix{} at $p=6$ to \numRitzValGapNine{} at $p=9$, against \numKdFreeMedMin{} to
\numKdFreeMedMax{} for the label-free transformer's three seeds, and at $p=9$ its seed means
were better than the label-free transformer's on every metric but the peak stress error. Its
median peak stress error, \numRitzValPeakSix{} at $p=6$ and \numRitzValPeakNine{} at $p=9$,
was \numRitzValPeakRatioMin{} to \numRitzValPeakRatioMax{} times the label-free
transformer's, and its peak stress error was the larger on at least \numRitzValFreePeakMin{}
of the instances at every degree. The energy-norm-supervised transformer had the lower
relative energy gap on at least \numRitzValKnormGapMin{} of the instances and the lower peak
stress error on all of them, at every degree, and lower seed means than the projection at
$p=9$ in relative energy gap, von Mises error and peak stress error. On the held-out
instances the comparisons read much the same: the energy-norm-supervised transformer had the
lower relative energy gap on at least \numRitzHoldKnormGapMin{} of them and the lower peak
stress error on all but at most \numRitzHoldKnormPeakMiss{}, at every degree. In two
dimensions, with \numRitzTwoKSix{} and \numRitzTwoKEight{} unknowns, the projection's
displacement error was smaller than the label-free transformer's on \numRitzTwoFreeDispMin{}
to \numRitzTwoFreeDispMax{} of the instances, depending on the seed and the degree, its
relative energy gap smaller on \numRitzTwoFreeGapSixMin{} to \numRitzTwoFreeGapSixMax{} at
$p=6$ and on \numRitzTwoFreeGapEightMin{} to \numRitzTwoFreeGapEightMax{} at $p=8$, and its
peak stress error larger on at least \numRitzTwoFreePeakMin{} of them. Supplementary
Section~\SuppRitz{} gives the shares for every network, set and degree. A smooth polynomial
space of a few hundred unknowns thus approximates the displacement field of these solids
better than each network, and their largest stresses worse than the transformers trained in
the energy norm.

\subsection{Inference and solving}\label{sec:cost:inference}

% generated from the committed records; do not edit
\begin{table}[tbp]
\centering
\footnotesize
\setlength{\tabcolsep}{4pt}
\caption{Wall time per instance of the label-free transformer's inference, on one GPU, and of solves of the same systems on the same machine's CPU, all load cases together and without the assembly, with the surrogate's accuracy on the same instances. CG: unpreconditioned conjugate gradients; its iterations (from the zero field / to the prediction's accuracy) and its displacement error are medians over the load cases, the other times and errors medians over the instances. Protocol in supplementary Section~\SuppTiming{}.}
\label{tab:cost}
\resizebox{\ifdim\width>\textwidth\textwidth\else\width\fi}{!}{%
\begin{tabular}{lccc}
\toprule
 & 2D & 3D & 3D, fine mesh \\
\midrule
Instances (load cases) & 32 (128) & 32 (128) & 8 (32) \\
Nodes, median & 562 & 4,960 & 38,156 \\
\midrule
Surrogate, first call & 3.2 ms & 49 ms & 2.15 s \\
Surrogate, repeated call & 2.4 ms & 49 ms & 2.16 s \\
\midrule
Sparse direct solve & 3.8 ms & 1.22 s & 133 s \\
CG from the zero field & 28 ms & 1.72 s & 26.1 s \\
CG from the prediction & 29 ms & 1.71 s & 27.4 s \\
CG to the prediction's accuracy & 8.8 ms & 292 ms & 4.50 s \\
the same, after rescaling by $c^*$ & 8.7 ms & 296 ms & 5.01 s \\
\midrule
Assembly of $K$, the repository's assembler & 9.5 ms & 900 ms & 8.07 s \\
Ritz projection after assembly & 1.4--2.1 ms & 34--144 ms & 0.54--1.39 s \\
\midrule
CG iterations per load case & 290 / 85 & 746 / 134 & 1,456 / 241 \\
\midrule
Surrogate's relative energy gap & 0.029 & 0.0085 & 0.12 \\
Surrogate's displacement error & 0.050 & 0.031 & 0.27 \\
Energy gap after rescaling by $c^*$ & 0.028 & 0.0077 & 0.065 \\
Displacement error after rescaling by $c^*$ & 0.040 & 0.014 & 0.076 \\
CG's displacement error at that accuracy & 0.049 & 0.017 & 0.16 \\
\bottomrule
\end{tabular}}

\smallskip
\begin{minipage}{\linewidth}\footnotesize\raggedright
Fine mesh, direct solve: the median over the seven of the eight instances that it solved before its budget of 900 s ran out.\par
CG to the prediction's accuracy reads the reference solution, so it is a lower bound; on the fine mesh it needs no iteration on the three load cases whose relative energy gap exceeds 1.
\end{minipage}
\end{table}


\Cref{tab:cost} times the label-free transformer's inference against solves of the same
systems on the same instances, with the sparse direct solver that produced the labels and
with conjugate gradients; the network ran on the GPU and the solvers on the CPU of the same
machine, and the assembly of the stiffness matrix, which the solvers and the Ritz projection
need and the surrogate does not, is listed apart (supplementary Section~\SuppTiming{} gives
the protocol). In three dimensions the direct solve took
a median \numCostThreeValDirectRatio{} times as long as the surrogate's first call on the
validation instances, per instance, and \numCostThreeFineDirectRatio{} times on the
\numCostFineDirectSolved{} fine-mesh instances that it solved before its time budget ran
out, where
the surrogate's median relative energy gap was \numCostThreeFineGap{} against
\numCostThreeValGap{} on the validation instances; in two dimensions it took
\numCostTwoDirectRatio{} times as long, and its solution is exact. Conjugate gradients from
the zero field, stopped at the first iterate as accurate as the prediction in the energy
norm, took a median \numCostThreeValMatchRatio{} and \numCostThreeFineMatchRatio{} times the
surrogate's first call on the three-dimensional validation and fine-mesh instances; this
stop reads the reference solution, so a practical run would take at least as long. The Ritz
projection took \numCostRitzValMin{} to \numCostRitzValMax{} per validation instance from
$p=6$ to $p=9$, against \numCostSurrVal{} for the surrogate's first call, after an assembly
that took \numCostAsmOverRitzMin{} to \numCostAsmOverRitzMax{} times as long as the
projection itself with the repository's assembler. We do not read these times as a speed
advantage of the surrogate: the network ran on a GPU and the solvers on a CPU, with
\numCostThreads{} threads (supplementary Section~\SuppTiming{}), neither the general sparse
LU factorisation of SuperLU nor unpreconditioned conjugate gradients is the solver of choice
for these symmetric positive definite systems, and the ratio of the matched solve to the
surrogate fell from \numCostThreeValMatchRatio{} to \numCostThreeFineMatchRatio{} on the
finer mesh.

\subsection{Labels and training}\label{sec:cost:labels}

In three dimensions the comparison consumed \numThreeSolves{} solves: \numThreeTrainSolves{}
for the 1,024 labelled training instances, \numThreeFewSolves{} for 64 fine-mesh instances
reserved for few-shot training and \numThreeEvalSolves{} for the in-band and fine-mesh
validation sets; in two dimensions it used \numCmSolves{}, \numCmSolvesTrain{} of them for
training. The validation solves are needed to evaluate any method; the training solves are
the price of supervision, and the label-free transformer needed none of them. They are cheap
here. At the median times of \cref{tab:cost}, the \numThreeTrainSolves{} three-dimensional
training solves take about \numCostLabelMin{} min of the direct solver, against
\numTimeThreeSeed{} h of training per seed on one NVIDIA GeForce RTX 5090, for 204,800
steps, which the displacement-supervised transformer needs as well. For these problems,
label-free training saves those minutes, not the cost of training. Full self-attention over
all nodes dominates that cost and grows with the square of the number of nodes: one training
step took \numEtwoTransformerFineStep{} s on a fine-mesh instance of \numEtwoNodesFine{}
nodes. A variant that attends over pooled tokens trained \numTimeBottleneckSpeedup{} times
faster but was less accurate, and was dropped under its pre-registered criterion
(supplementary Section~\SuppSide{}). Set against the direct solve, the training of one seed
is recovered after about \numCostBreakEvenDirect{} predicted instances, and after about
\numCostBreakEvenMatch{} against conjugate gradients stopped at the surrogate's accuracy, an
upper bound since that stop reads the reference solution. Label-free training pays for
itself where a solve is expensive enough (\cref{sec:discussion}). As warm starts the
predictions did not shorten conjugate gradients, here or in the pre-registered test of the
July run (K5 in \cref{tab:claims}); supplementary Section~\SuppWarm{} gives these
measurements, the iteration bound that explains them and a comparison with reported
speed-ups.
